\uuid{33gw}
\exo7id{3164}
\titre{exo7 3164}
\auteur{quercia}
\organisation{exo7}
\datecreate{2010-03-08}
\isIndication{false}
\isCorrection{false}
\chapitre{Polynôme, fraction rationnelle}
\sousChapitre{Autre}
\cpge{mpsi-mp2i.espaces-vectoriels.c, pcsi.espaces-vectoriels.c, ptsi.espaces-vectoriels.c, mpsi-mp2i.polynomes-fractions-rationnelles.c, pcsi.polynomes.c, ptsi.polynomes.c}
\module{Algèbre}
\niveau{L1}
\difficulte{}

\contenu{
\question{
Soit $a,b \in  K,\ a \ne b$. On pose $P_k = (X-a)^k(X-b)^{n-k}$.
D{\'e}montrer que la famille $(P_0,\dots,P_n)$ est libre.
}
}
